\subsection{The separable closure of a field}

DEFINITION 4. --- A field K is said to be separably closed if every separable algebraic extension of K is trivial.

An algebraically closed field is separably closed. Conversely, if a perfect field K is separably closed, it is algebraically closed, because every algebraic extension of K is separable (V, p.~3, Prop. 2).

DEFINITION 5. --- Let K be a field. By a separable algebraic closure, or (by abuse of language) separable closure of K we understand any extension E of K which is algebraic and separable over K, and such that the field E is separably closed.

When K is perfect, there is complete identity between the notions of separable closure and algebraic closure of K (V, p.~36, Prop. 2 and p.~43, Cor. 1).

PROPOSITION 14. --- Let \(\Omega\) be an algebraically closed extension of K.

\begin{enumerate}
\def\labelenumi{\alph{enumi})}
\item
  The relative separable algebraic closure \(\Omega_s\) of \(\mathbf{K}\) in \(\Omega\) is a separable closure of \(\mathbf{K}\) .
\item
  If E and E' are two separable closures of K, there exists a K-isomorphism of E onto E'.
\end{enumerate}

Let \(F\) be a separable algebraic extension of \(\Omega_s\) ; since \(\Omega\) is algebraically closed, there exists an \(\Omega_s\) -homomorphism \(u\) of \(F\) into \(\Omega\) (V, p.~20, Th. 1). By Prop. 9 (V, p.~42) \(u(F)\) is separable over \(K\) , hence \(u(F) = \Omega_s\) . Therefore \(F\) is a trivial extension of \(\Omega_s\) and so \(\Omega_s\) is separably closed, whence \(a\) ).

Let E be a separable closure of K. Since E is an algebraic extension of K, there exists a K-homomorphism v of E into \(\Omega\) (V, p.~20, Th. 1). Hence \(v(\mathrm{E})\) is separable algebraic over K, whence \(v(\mathrm{E}) \subset \Omega_{s}\) . By V, p.~42, Prop. 9, \(\Omega_{s}\) is separable over \(v(\mathrm{E})\) and since the field \(v(\mathrm{E})\) is separably closed, we have \(v(\mathrm{E}) = \Omega_{s}\) . It follows that v is a K-isomorphism of E onto \(\Omega_{s}\) . Now b) is an immediate consequence.

COROLLARY. --- Let E be a separably closed extension of K and F a separable algebraic extension of K; then there exists a K-homomorphism of F into E.

Let \(\Omega\) be an algebraic closure of E; we have \(\Omega_s \subset E\) and it suffices to treat the case where \(E = \Omega_s\) . Since F is an algebraic extension of K, there exists a K-homomorphism \(u\) of F into \(\Omega\) (V, p.~20, Th. 1). Since the field \(u(F)\) is separable over K, we have \(u(F) \subset \Omega_s\) and \(u\) defines a K-homomorphism of F into \(\Omega_s = E\) .

\pandocbounded{\includegraphics[keepaspectratio]{images/part001_42ca6df9.jpg}}

Remarks. --- 1) Let E and E' be two separable closures of K. If K is not separably closed, there exist several K-isomorphisms of E onto E'. For E is then a non-trivial Galois extension of K, and so there exist K-automorphisms of E distinct from the identity (V, p.~56, Th. 1).

2) Let \(E\) be an algebraic and separable extension of \(K\) . If every algebraic and separable extension of \(K\) is isomorphic to a subextension of \(E\) , then \(E\) is a separable closure of K. For if \(E'\) is a separable closure of K, then each of the extensions E and \(E'\) is isomorphic to a subextension of the other; hence E and \(E'\) are isomorphic extensions of K (V, p.~52, Prop. 1, a)).

